%==================================================================
% Ini adalah abstrak dalam bahasa inggris 
%==================================================================

%% DILARANG EDIT BAGIAN INI
\clearpage
\phantomsection
\addcontentsline{toc}{chapter}{ABSTRACT}
\begin{center}
    \textbf{\large{\judulen}}\\[2em]
    \penulis\\
    NIM \nim\\[2em]
    \textbf{\textit{ABSTRACT}}\\[0.5cm]
\end{center}
%% DILARANG EDIT BAGIAN INI

%% edit bagian ini
\textit{INFINITE is a division under the Rekayasa Teknologi (Restek) Student Activity Unit at Universitas Negeri Yogyakarta (UNY) that serves as a platform for student self-development in the field of information technology. The previous INFINITE organizational information system was fragmented into two systems, namely Laravel with a monolithic architecture on shared hosting and a separate Strapi CMS, creating challenges in maintenance, data consistency, operational cost inefficiency, and the inability to integrate with external applications and the Kubernetes infrastructure. This study aims to develop a containerized INFINITE organizational information system with a headless architecture using Laravel as the backend, Next.js as the frontend, and authentication based on the OpenID Connect (OIDC) protocol, deployed on Kubernetes using the GitOps approach with ArgoCD.}

\textit{System development was carried out using the Research and Development (R\&D) method with the Waterfall model through four stages: Requirements Definition, System and Software Design, Implementation and Unit Testing, and Integration and System Testing. Testing was conducted based on the ISO/IEC 25010:2023 standard for \textit{Product Quality} and ISO/IEC 25019:2023 for \textit{Quality in Use}, covering the aspects of Functional Suitability, Interaction Capability, Performance Efficiency, and Beneficialness.}

\textit{The results show that the system was successfully developed and deployed as a containerized application on Kubernetes, with testing results of Functional Suitability at 100\% (Very Feasible), Interaction Capability at 91.73\% (Very Feasible), Performance Efficiency on the frontend with an average Lighthouse score of 90.42 (Good) and on the backend with a Response Time of 1 second, Throughput of 53.6 requests per second, Error Rate of 0\%, and a maximum of 189 simultaneous Virtual Users, and Beneficialness at 81.30\% (Very Feasible). Based on these results, the INFINITE organizational information system is considered feasible as a solution for integrated and sustainable organizational data management.}\\[0.6cm]
%% edit sampai sini

%% DILARANG EDIT BAGIAN INI
\noindent \textit{\textbf{Key words:} Information System, Headless, Container.}
%% DILARANG EDIT BAGIAN INI